\section{Kesimpulan}

\begin{enumerate}
	\item Ketiga algoritma \textit{local search} (Hill-Climbing dengan keempat variannya,
	Simulated Annealing, dan Genetic Algorithm) berhasil diimplementasikan untuk
	persoalan pengisian peti kemas truk. Seluruh solusi akhir memenuhi batasan dinding,
	tumpang tindih, tumpuan, \textit{fragile}, dan kapasitas berkat heuristik
	\textit{repair}.
	\item Simulated Annealing memberikan keseimbangan terbaik antara kualitas dan waktu.
	SA mencapai optimum global (terbukti oleh batas atas \textit{knapsack}) pada instans
	\textit{small} di ketiga \textit{run}, dan berada 5--10\% dari batas atas pada
	instans yang lebih besar dalam hitungan detik.
	\item HC \textit{random-restart} memberikan hasil terbaik dan paling konsisten pada
	instans \textit{medium} dan \textit{large}, tetapi waktunya paling lama (sekitar 28
	menit per \textit{run} pada \textit{large}) karena membangkitkan seluruh tetangga di
	setiap iterasi. HC \textit{stochastic} paling cepat tetapi paling mudah terjebak.
	\item Genetic Algorithm versi salindia (\textit{roulette wheel}, \textit{single-point
	crossover}, mutasi) bekerja cukup baik pada instans kecil, tetapi tertinggal jauh pada
	instans sedang dan besar. \textit{Crossover} yang menukar ekor kromosom merusak susunan
	geometris yang dibangun induk. Menambah populasi dan generasi tetap meningkatkan
	hasil, tetapi tidak menutup selisih tersebut.
	\item Parameter terpenting SA adalah banyak iterasi dan jadwal suhu. Jadwal geometrik
	dengan $p_0$ dan $T_{\min}$ kecil memberikan hasil terbaik.
	\item Objektif alternatif berbasis ETA berhasil menggeser pilihan paket ke paket yang
	lebih mendesak dengan pengorbanan total \textit{value} yang kecil (3,8\%).
	\item Dukungan multi-truk diperoleh hanya dengan menambahkan identitas truk pada
	\textit{state} serta memeriksa batasan per truk, tanpa mengubah ketiga algoritma.
	Dengan dua truk, SA dan HC \textit{steepest} mengangkut 14--19\% lebih banyak
	\textit{value} dan berada 2--4\% dari batas atas.
\end{enumerate}

\section{Saran}

\begin{enumerate}
	\item \textbf{Representasi GA berbasis urutan.} Kromosom dapat diubah menjadi urutan
	paket beserta orientasinya, lalu didekode secara \textit{greedy} ke \textit{extreme
	point} pertama yang valid. Dengan \textit{order crossover}, anak selalu berupa susunan
	valid yang mewarisi urutan prioritas induknya.
	\item \textbf{Inisialisasi dan pengisian \textit{greedy}.} Uji coba awal di luar kode
	utama, yaitu menambahkan langkah pengisian \textit{greedy} ke \textit{extreme point}
	setelah \textit{repair}, menaikkan rata-rata GA (populasi 50, generasi 100) dari
	sekitar 1439 menjadi 2313 pada \textit{medium} dan dari 2302 menjadi 5370 pada
	\textit{large}. Sebagian besar kenaikan berasal dari heuristik \textit{greedy} itu
	sendiri, sehingga pendekatan ini perlu dikaji sebagai algoritma memetik.
	\item \textbf{Evaluasi inkremental.} \textit{Repair} dan pemeriksaan batasan saat ini
	menghitung ulang seluruh \textit{state} untuk setiap tetangga. Evaluasi inkremental
	(hanya memeriksa paket yang berubah) akan mempercepat HC \textit{steepest} secara
	signifikan.
	\item \textbf{Eksekusi paralel.} \textit{Run} antar-\textit{seed} dan antarkonfigurasi
	independen, sehingga eksperimen dapat dijalankan paralel. Pengukuran durasi
	sebaiknya tetap dilakukan dalam kondisi yang sama untuk semua algoritma.
\end{enumerate}
